\section{A comparison map}
\label{app:patterns}

The table condenses the comparison developed in the article. Its right column
describes common agent defaults, not necessary properties of generated text.
The same patterns occur in human drafts, which is why the comparison belongs
at the level of passages rather than authors.

\begingroup
\small
\setlength{\tabcolsep}{4pt}
\begin{longtable}{@{}>{\raggedright\arraybackslash}p{0.16\textwidth}
  >{\raggedright\arraybackslash}p{0.39\textwidth}
  >{\raggedright\arraybackslash}p{0.39\textwidth}@{}}
\toprule
\textbf{Dimension} & \textbf{Strong human prose} &
\textbf{Common agent default} \\
\midrule
\endfirsthead
\toprule
\textbf{Dimension} & \textbf{Strong human prose} &
\textbf{Common agent default} \\
\midrule
\endhead

Reader model &
Chooses background and examples for what a particular reader knows and needs. &
Falls back on portable explanation and generic helpfulness when the reader is
unspecified. \\
\addlinespace

Priority &
Gives more space to the central claim and deciding evidence than to routine
background. &
Distributes equal visual and rhythmic weight across background, evidence,
caveats, and implications. \\
\addlinespace

Semantic fit &
Matches subjects with actions they can perform and names missing links such as
extent, applicability, or consequence. &
Produces grammatical but loose combinations such as claims that ``stop'' or
sections that ``know.'' \\
\addlinespace

Evidence &
Names sources and distinguishes direct results from interpretation or open
questions. &
Relies on vague attribution, inflated significance, or evidential verbs
stronger than the available support. \\
\addlinespace

Continuity &
Connects sentences through cause, condition, contrast, and consequence. &
Places related ideas in separate, similarly shaped sentences or list items and
leaves the reader to recover the dependency. \\
\addlinespace

Cadence &
Varies sentence length and structure because ideas have different roles and
weights. &
Stacks clipped sentences for drama or repeats medium-length constructions with
the same opening and closure. \\
\addlinespace

Contrast &
Chooses syntax according to the comparator, negative result, mechanism, or
claim boundary at issue. &
Repeats frames such as ``not X but Y'' or ``treats X as Y, not Z'' after the
underlying relation changes. \\
\addlinespace

Voice &
Reveals judgment through selected details, calibrated uncertainty, and
deliberate irregularity. &
Defaults to a portable middle register assembled from corporate, social-media,
documentation, and chat conventions. \\
\addlinespace

Structure &
Uses headings and lists when readers need to navigate a real change of question
or scan a complete set. &
Uses visible organization to make a document look complete before the
dependencies among its ideas are clear. \\
\addlinespace

Abstract &
Explains the problem, central answer, decisive evidence, and claim boundary. &
Tours the sections, inventories contributions, or gives every component equal
space. \\
\addlinespace

Ending &
Ends on the supported consequence, limitation, unresolved tension, or
concrete fact at which the reasoning arrives. &
Appends a summary, slogan, broad implication, or optimistic future because the
form appears to require closure. \\

\bottomrule
\end{longtable}
\endgroup
